\documentclass[10pt,a4paper]{amsart}
\usepackage[margin=19mm]{geometry}
\usepackage{amsmath,amssymb,mathtools}
\usepackage[scr=rsfs,cal=euler]{mathalfa}
\usepackage{xfrac}
\usepackage[unicode,hidelinks,bookmarks=false]{hyperref}
\newcommand{\ie}{i.e.\ }
\newcommand{\cf}{cf.\ }
\newcommand{\resp}{respectively\ }
\newcommand{\Subsection}{\subsection}
\providecommand{\nobreakdash}{\nobreak\hbox{-}}
\setlength{\emergencystretch}{2em}
\multlinegap0pt
\usepackage{bm}
\usepackage{bbm}
\usepackage{dsfont}
\usepackage{xspace}
\usepackage{cancel}
\usepackage{mathtools}
\usepackage{amsthm}
\makeatletter
\@ifundefined{lemma}{\newtheorem{lemma}{Lemma}}{}
\makeatother
\newcommand{\Span}{\mathrm{Span}}
\newcommand{\braces}[1]{\left\{ #1 \right\}}
\newcommand{\cH}{\mathcal{H}}
\newcommand{\cM}{\mathcal{M}}
\newcommand{\cV}{\mathcal{V}}
\newcommand{\mA}{\mathbf{A}}
\title[Decentralized Optimization with Coupled Constraints]{Decentralized Optimization with Coupled Constraints}
\author{Demyan Yarmoshik, Alexander Rogozin, Nikita Kiselev, Daniil Dorin, Alexander Gasnikov, Dmitry Kovalev}
\date{Source: arXiv:2407.02020v4 (CC BY 4.0)}
\begin{document}
\maketitle

\noindent\emph{Source and licence.} Decentralized Optimization with Coupled Constraints. Version arXiv:2407.02020v4 (CC BY 4.0), distributed under the Creative Commons Attribution 4.0 International licence, as recorded in the arXiv licence metadata for this exact version and shown on the abstract page https://arxiv.org/abs/2407.02020v4. Full rights notice: licenses/arxiv-2407.02020-CC-BY-4.0.txt.\par\smallskip

\noindent\textit{Editorial scope.} Counted unit: the Lemma whose source label is \texttt{eq:nonzero\_components\_bound} in the article (source0.tex lines 1284--1304, source label \texttt{eq:nonzero\_components\_bound}), delivered together with the article's own complete original proof of it (source0.tex lines 1305--1364, one \texttt{proof} environment of 60 lines). No mathematics is written, repaired or re-derived: statement and proof are the article's own text line for line. Required context retained uncounted: none. Every definition, display and label of those lines is retained unchanged. Omitted from this package and not redistributed: 1--1283, 1365--1473. Nothing inside a retained excerpt is omitted or altered: every retained body is delivered exactly as the article's lines are written, so no omission receipt of any kind accompanies this package. Citations: the retained excerpts carry no citation command at all, so no bibliography entry of the article is transcribed and no reference is redistributed. Reference adaptation: none was needed and none was made; every reference inside the delivered unit resolves inside it, so no unresolved reference or citation is produced. Printed slips kept literal: no printed word, symbol, hypothesis or number of the counted statement or of its proof was altered. Layout and class: the article's own class is \texttt{article} with packages including iclr2025\_conference, times, amsthm, amssymb, amsmath, inputenc, fontenc, hyperref, url, booktabs, nicefrac, microtype, wrapfig, lipsum; the wrapper is the lane's pinned \texttt{amsart} editorial wrapper at 10pt on A4, which restates the theorem-like environments the retained text needs and adds the article's own macro definitions that the retained text needs (\texttt{\textbackslash Span}, \texttt{\textbackslash braces}, \texttt{\textbackslash cH}, \texttt{\textbackslash cM}, \texttt{\textbackslash cV}, \texttt{\textbackslash mA}), with extra wrapper packages (bm, bbm, dsfont, xspace, cancel, mathtools). The delivered numbering is the unit's own. Assets: no retained span contains a figure environment, a graphics-inclusion command, a PDF-page inclusion, a table or a TikZ picture; no image file is copied into this package, the package declares no asset dependency and no image file is redistributed with it. Licence: version arXiv:2407.02020v4 of this article is distributed under the Creative Commons Attribution 4.0 International licence, as recorded in the arXiv licence metadata for this exact version and shown on the abstract page https://arxiv.org/abs/2407.02020v4. A full rights notice is delivered at licenses/arxiv-2407.02020-CC-BY-4.0.txt. Characters: every retained excerpt is pure ASCII, so the delivered engine, XeLaTeX with its default Latin Modern text font, reports no missing character.
\begin{lemma}
	Let $s_i(k)$ denote the maximum index of a nonzero component of vector blocks $p, t$ held by $i$-th node at step $k$ in its local memory $\cM_i(k)$ and vector $z$ held by $i$-th node in its local memory $\cH_i(k)$, i.e.
	\begin{align*}
		s_i(k) = 
		\begin{cases}
			0, &\hspace{-1.0cm}\text{if } \cM_i(k)\subseteq\braces{0} \text{ and } \cH_i(k)\subseteq\braces{0} \\
			\min\Big\{s\in \braces{1, 2, \ldots}: \cH_i(k)\subseteq\Span\braces{e_1, \ldots, e_s} \\ \qquad~ \text{ and } \cM_i(k)\subseteq \Span\braces{e_1, \ldots, e_s} \times\braces{e_1, \ldots, e_s}\Big\}, &\text{else}.
		\end{cases}
	\end{align*}
%	\begin{align*}
%		v_i(k) = 
%		\begin{cases}
%			0, &\cH_i(k)\subseteq\braces{0} \\
%			\min\braces{v\in \braces{1, 2, \ldots}: \cH_i(k)\subseteq \Span\braces{e_1, \ldots, e_v}}, &\text{else}.
%		\end{cases}
%	\end{align*}
	Let $k_q$ denote the number of algorithm step by which exactly $q$ communication steps have been performed, where $q\geq 0$. For any $k\in\braces{1, \ldots, k_q}$ we have
	\begin{align}\label{eq:nonzero_components_bound}
		\max_i s_i(k)\leq 2 + \left\lfloor \frac{q}{\frac{n}{3}+ 1} \right\rfloor
	\end{align}
\end{lemma}

\begin{proof}
	If the method performs a multiplication by $\mA_i$, it transfers the information from $\cM_i(k)$ to $\cH_i(k)$. If multiplication by $\mA_i^\top$ is performed, the information is transferred in the opposite direction, i.e. from $\cH_i(k)$ to $\cM_i(k)$. If the method performs a matrix multiplication step (either by $\mA$ or $\mA^\top$), then from the structure of $\mA_i$ we obtain
	\begin{align*}
		s_i(k+1)\leq s_i(k) + \begin{cases} 1 - (s_i(k) \mod 2), &i\in\cV_1 \\ 0, &i\in\cV_2 \\ (s_i(k)\mod 2), &i\in\cV_3 \end{cases}
	\end{align*}
%	If a method performs a multiplication by $\mA_i^\top$, it transfers information from $\cH_i(k)$ to $\cM_i(k)$ and we obtain
%	\begin{align*}
%		r_i(k+1)\leq \begin{cases} \max(r_i(k), s_i(k) - (s_i(k)\mod 2)) + \mathbb{I}(s_i(k) = 1), &i\in\cV_1 \\ r_i(k), &i\in\cV_2 \\
%		\max(r_i(k), s_i(k) - ((s_i(k) + 1) \mod 2), &i\in\cV_3
%		\end{cases}
%	\end{align*}
%	If the method makes a communication round, it transfers the information between $\cH_i(k)$ for different $i$. Due to the structure of network, if the method makes a communication round, it follows
%	\begin{align}\label{eq:s_update_comm}
%		s_i(k+1)\leq 
%		\begin{cases}
%			\max\cbraces{s_{i-1}(k), s_i(k), s_{i+1}(k)}, &i\in \braces{2, \ldots, n - 1} \\
%			\max\cbraces{s_1(k), s_2(k)}, &i = 1 \\
%			\max\cbraces{s_{n-1}(k), s_n(k)}, &i = n
%		\end{cases}
%	\end{align}
	We will prove that\\
	1. For $q = 2\ell(n/3 + 1),~ \ell\in\braces{0, 1, \ldots}$ we have
	\begin{align}\label{eq:info_flow_even}
		s_i(k_q)\leq
		\begin{cases}
			1 + 2\ell,~~ i\in\cV_1 \\
			1 + 2\ell,~~ i\in\cV_2 \\
			2 + 2\ell,~~ i\in\cV_3
		\end{cases}
%		s_{n/3}(k_q), s_{n/3+1}(k_q), \ldots, s_{2n/3+1}(k_q) \leq 1 + 2\ell,~~~~ s_{2n/3+2}(k_q) \leq 2 + 2\ell.
	\end{align}
	2. For $q = (2\ell + 1)(n/3 + 1),~ \ell\in\braces{0, 1, \ldots}$ we have
	\begin{align}\label{eq:info_flow_odd}
		s_i(k_q)\leq
		\begin{cases}
			2 + (2\ell + 1),~~ i\in\cV_1 \\
			1 + (2\ell + 1),~~ i\in\cV_2 \\
			1 + (2\ell + 1),~~ i\in\cV_3
		\end{cases}
%		s_{n/3}(k_q) \leq 2 + (2\ell + 1),~~~~ s_{n/3+1}(k_q), \ldots, s_{2n/3+1}(k_q), s_{2n/3+2}(k_q) \leq 1 + (2\ell + 1).
	\end{align}
	The proof follows by induction.
	
	\noindent\textbf{Induction basis}. Let $q = 0$. From definitions of $f_i$ and $\mA_i$ it follows that
	\begin{align*}
			s_i(k_0) \leq 
			\begin{cases}
			1,~~ i\in\cV_1 \\
			0,~~ i\in\cV_2 \\
			2,~~ i\in\cV_3
		\end{cases}
	\end{align*}
	Therefore, for $q = 0$ our statement holds.
	
	\noindent\textbf{Induction step for $q = (2\ell + 1)(n/3 + 1)$}. Consider $q_- = q - n/3 = 2\ell(n/3 + 1)$. From \eqref{eq:info_flow_even} we have that for the spread of nonzero components from $\cV_3$ to $\cV_1$ it requires $n/3$ communication rounds to reach node $n/3 + 1$. After one more communication round, the information reaches node $n/3$.
	
	\noindent\textbf{Induction step for $q = 2\ell (n/3 + 1)$}. The proof follows by the same argument as for $q = (2\ell + 1)(n/3 + 1)$.

	\vspace{0.5cm} We just proved the statement of lemma, i.e. relation \eqref{eq:nonzero_components_bound}, for $q$ divisible by $(n/3 + 1)$. Between such checkpoints, the information (i.e. the number of nonzero components) traverses nodes of $\cV_2$ and therefore $\max_i s_i(k)$ stays unchanged. Thus the statement of the lemma is proven. 
\end{proof}

\end{document}
